\section{From moments to logarithmic length windows}\label{L:part11}

We prove Theorem~\ref{L:thm:window}. The first step is to retain a substantial length-weighted mass after cutting off exceptional diameters and lengths. The second step uses the length-controlled growth just proved to attach two disjoint connectors to the surviving arches and controls the number of ways to recover the input from an output bridge. We keep \(\mu(\omega)=X^{|\omega|}\); counts without a spatial root on plane polygons are modulo the primitive translations.

We use the constant-factor cylinder estimate \eqref{L:eq1} and its plane ball upper bound \(C r^{1/12}\) at diameter scale \(r\), together with the following strip consequence. In an \emph{infinite strip} between horizontal lattice-side walls at distance \(h\), for a point \(v\) in the middle third (vertex or face center) the enclosing gas satisfies
\[
 {\cal Z}_{\rm strip}(v)\ge c h^{1/12}(\log h)^{-1/12}.
\]
Indeed use \eqref{L:eq1} on the regular cylinder with even period \(n\) along the walls of order \(h/(C_1\log h)\), \(C_1\) sufficiently large. Winding separators can be dropped at cost a constant factor as before. The single-loop weight of contractible separators of height \(>h/6\) is negligible by the exponential pressure bound (multiply by height plus a constant to include placements at the slice). Removing all configurations containing one thus costs a negligible fraction of the contractible gas. All others lift around the corresponding two marks \emph{inside infinite strips} of the required height, regardless of their width, splitting into two plane gas spaces as before. This and the vertex transfer (bounded mass of polygons through the vertex) give the bound.

\subsection{Aspect cutoff}\label{L:step18}\label{L:prop:aspect}

Let $x,y$ be normal level coordinates in two distinct lattice-side
 directions.  Write $K_D(h)$ for the mass of strict $x$-bridges of height
$D$ from a fixed source whose absolute $y$-deviation from that source is
at most $h$.  We first prove
\begin{equation}\label{L:eq:aspect-bound}
 K_D(h)\le C D^{-1/4}e^{-cD/h}\qquad(D\ge h\gg1).
\end{equation}
The cylinder tail will supply the exponential factor.  To use it, we
join two long, narrow bridges into a polygon, retaining enough possible
cut heights that the construction can be repeated at constant cost.

Fix $D\ge h$, and choose an integral gap $g_0$ comparable to $10h$
on an $x$-wall.  A step along this wall changes $y$ by one.  Start two
bridges on that wall with gap in $[g_0-2h,g_0+2h]$.  For each
$d\in[D,2D]$, the two translated bridge increments have product mass
$K_d(h)^2$.  Interchanging the increments reverses the change in the
terminal gap.  At least one of these two choices leaves that gap in
the same interval, so the retained product mass is at least
$K_d(h)^2/2$.  The two paths are disjoint because their starting gap
exceeds $2h$.

Extract the first simultaneous qualifying $x$-cut in $[D,2D]$
whose prefixes have the required confinement and terminal gap, allowing
the terminal wall as a cut.  The terminal wall qualifies.  From a selected cut, the two discarded
suffixes are strict bridges of the same height.  Summing over possible
terminal walls bounds their total product mass by
$\sum_{j\le D}B_j^2\le C\sqrt D$, including the empty suffix at $j=0$.
Thus, for every admissible starting gap, the family of distinct
prefix pairs has mass at least the common quantity
\[
 W:=cD^{-1/2}\sum_{d=D}^{2D}K_d(h)^2.
\]
The first-cut convention recovers this pair after any continuation in
the next $x$-slab.

Repeat the construction for $k$ stages.  At each stage the starting
gap is admissible, so the same lower bound $W$ applies.  Close the two
initial ends and the two terminal ends by exterior caps, each of
mass at least $ch^{-5/4}$ and diameter $O(h)$.  The resulting polygon
has $x$-extent between $kD-O(h)$ and $2kD+O(h)$.  Although its lateral
center may drift between stages, its spread in $y$ at any fixed
$x$-level is $O(h)$: the stages occupy disjoint $x$-slabs and the
terminal gap is always comparable to $h$.  It therefore projects
without collision to the regular cylinder of circumference $C'h$,
for a sufficiently large fixed $C'$.

For fixed $h,D$, the number of inputs producing a cylinder polygon is
polynomial in $k$.  Choose a starting port on its lift and the terminal
level at which to remove the two caps; there are polynomially many
choices in $k,h,D$.  The successive pair cuts are then read by the
first-cut rule.  The cylinder height tail, with the deterministic
length bound used to absorb these recovery choices, consequently gives
\[
 c h^{-5/2}W^k
 \le \operatorname{poly}(k,h,D)\,h^C e^{-ckD/h}.
\]
Taking $k$th roots and then $k\to\infty$ yields $W\le C e^{-cD/h}$,
and hence
\[
 \sum_{d=D}^{2D}K_d(h)^2\le C\sqrt D\,e^{-cD/h}.
\]
To obtain one prescribed height, enclose the paths counted by $K_D(h)$
in their parallelogram and clip at an $x$-level $d$ in a proportional
bin below $D$.  Positive boundary flux gives
$K_D(h)\le C K_d(Ch)$ at each such level.  Sum its square over the
$\asymp D$ available levels and use the preceding bin bound.  This
proves \eqref{L:eq:aspect-bound} when $D/h$ is large; the ordinary
bound $K_D(h)\le B_D\le CD^{-1/4}$ covers bounded ratios.

It follows that arches and bridges in a $y$-strip of height comparable
to $h$, from one fixed wall source, satisfy
\[
 \mu\{\operatorname{diam}\omega>s\}
 \le C h^{-1/4}e^{-cs/h}\qquad(s\ge C_0h).
\]
Indeed clip by two $x$-walls at separation a small fixed multiple of
$s$, so every such path is lost.  Average over $\asymp s$ central
wall sources and reverse the new exits.  There are $O(h)$ ports on
the new walls, and reaching those central sources costs at most
$C K_{\lfloor c's\rfloor}(Ch)$ per reversed source.  The preceding
estimate gives $C(h/s)s^{-1/4}e^{-cs/h}$, which is at most the claimed
bound.  Exhaustion by larger parallelograms justifies the argument in
an infinite strip.

\subsection{Length moments}\label{L:step19}\label{L:prop:moments}

For diameter at most $s$, let $F_{A,p}(s,g)$ be the length-$p$ mass of
arches from a fixed source with gap in $[g,2g)$, and let
$F_{B,p}(s,d)$ be the length-$p$ mass of bridges from a fixed source
at exact height $d$.  With $b(s)=1+\log\log(e^e+s)$, we prove, for
$p\in\{1,3/2\}$ and $1\le g\lesssim s$,
\begin{equation}\label{L:eq:geometric-moments}
 F_{A,p}(s,g)\le Cb(s)s^{p/3}g^{p-1/4},\qquad
 \sum_{1\le d\le s}d^{-1/4}F_{B,p}(s,d)
 \le Cb(s)s^{1/2+4p/3}.
\end{equation}
Fixed changes between diameter and level units only change the constants.
Both estimates come from sewing a path into a polygon.  The main issue
is how many cuts of the resulting polygon can recover an input path.

A simple double cut is a lattice-side grid line meeting a polygon at
exactly two crossing ports.  For fixed comparability constants, let
$M_g(P)$ count such cuts in the required orientation, with crossing
gap comparable to $g$ and at distance at most $Cg$ above the polygon's
bottom.  For every fixed integer $k\ge1$,
\begin{equation}\label{L:eq:fixed-cut-moment}
 \sum_P\mu(P)M_g(P)^k\le C_k g^{-2+k/2},
\end{equation}
where polygons are counted modulo primitive translations.  To see
this, first select $r$ distinct cuts in increasing height order.
Translation fixes one crossing of the first cut, and its partner has
$O(g)$ choices.  The lower and upper outer arcs have specified ends
at gaps comparable to $g$, so each costs at most $Cg^{-5/4}$.
Between two consecutive cuts, polygon connectivity forces two strict
bridges; if their height difference is $j$, their product mass is at
most $B_j^2$.  Summing each positive gap up to $Cg$ costs $O(\sqrt g)$.
The total is therefore
\[
 C_r g\,g^{-5/2}(\sqrt g)^{r-1}=C_r g^{-2+r/2}.
\]
Expanding $M_g^k$ by its number $r\le k$ of distinct levels proves
\eqref{L:eq:fixed-cut-moment}.

For an arch at gap comparable to $g$, attach below its wall a cap of
mass at least $cg^{-5/4}$ and diameter $O(g)$.  Its original wall is
one of the cuts counted by $M_g$.  Choosing that cut and one crossing
recovers the arch and cap up to a fixed number of orientations.
Consequently
\[
 c g^{-5/4}F_{A,p}(s,g)
 \le C\sum_{\operatorname{diam}P\le Cs}\mu(P)|P|^p M_g(P).
\]
Let $Q_2(s)=\sum_{\operatorname{diam}P\le Cs}\mu(P)|P|^2$.
Summing \eqref{L:eq10} over dyadic scales gives
$Q_2(s)\le Cb(s)s^{2/3}$.  H\"older's inequality with
$k=2/(2-p)$ bounds the last display by
\[
 C Q_2(s)^{p/2}
 \left(\sum_P\mu(P)M_g(P)^k\right)^{1/k}.
\]
After multiplying by $g^{5/4}$, the power of $g$ is
$5/4+(-2+k/2)/k=p-1/4$, giving the first estimate.

For a bridge of height $d\le s$, choose a second bridge of the same
height, of diameter $O(d)$ and mass at least $cd^{-1/4}$ by transverse
tightness.  Translate it by one of $\asymp s$ shifts at distance
comparable to a sufficiently large fixed multiple of $s$.  The two
bridges are disjoint; caps of gaps comparable to $s$ close them at
both walls.  The two caps contribute $cs^{-5/2}$, and the shift
choices contribute $cs$, for a net factor $cs^{-3/2}$ besides the
second-bridge mass.  The output has diameter $O(s)$.  Its two joining
walls are counted by $M_s$, so the recovery multiplicity is
$O(M_s^2)$.  Thus
\[
 c s^{-3/2}\sum_{1\le d\le s}d^{-1/4}F_{B,p}(s,d)
 \le C\sum_{\operatorname{diam}P\le Cs}\mu(P)|P|^pM_s(P)^2.
\]
Use H\"older with $k=4/(2-p)$ and power $2/k$ on the cut moment.
The resulting power of $s$ is
$3/2+p/3+(-2+k/2)2/k=1/2+4p/3$.
Since $b(s)^{p/2}\le b(s)$, this proves the second estimate.

\subsection{Proof of the logarithmic window}\label{L:step20}
\begin{proof}[Proof of Theorem~\ref{L:thm:window}]
Put
\[
 s=h(\log h)^{1/128},\qquad R_0=h(\log h)^{1/64},\qquad
 \ell_-=h^{4/3}(\log h)^{-1/8},\qquad
 \ell_+=h^{4/3}(\log h)^{1/2}.
\]
We first retain paths with lengths in $[\ell_-,\ell_+]$ and diameter
at most $s$.  These paths are either bridges already or arches that
will be completed outside their supporting half-plane.  The final
recovery count will determine the cost of this completion.

\medskip\noindent\textbf{Retained path mass.}
For each integer $d\in[h,2h]$, call a chord in the horizontal strip
of height $d$ deep if it visits the middle third.  The strip enclosing
gas lower bound and \eqref{L:eq6} give length-weighted mass at least
\[
 c h^{13/12}(\log h)^{-1/12}
\]
per fixed wall source.  To justify the strip limit, exhaust by long
convex parallelograms.  Chords from a remote wall visiting a fixed
middle vertex disappear: reverse from that wall and clip before
reaching the vertex, then apply the aspect cutoff to its $O(h)$
possible source ports.  Sum the resulting infinite-strip comparison
over the $\asymp h$ middle vertices representing translation classes
along the walls.  Translation turns the sum over all chord endpoints
into middle visits by paths from one representative source on each
wall; the two wall sums agree by symmetry.  This proves the displayed
lower bound.  The unweighted mass of deep paths is at most
$Ch^{-1/4}$, by the bridge bound and the arch height tail.

Let $\bar\mu$ be the average of these deep-path measures over the
$\asymp h$ heights in $[h,2h]$.
The length-$3/2$ estimate gives, for diameter at most $t\ge h$,
\[
 Cb(t)(t/h)^{5/2}h^{7/4}.
\]
For arches this follows by summing the dyadic gap estimates up to
$t$; for bridges, remove the factor $d^{-1/4}$ at $d\asymp h$
and divide the height sum by $h$.  Fixed enlargements of diameter
cutoffs absorb the comparison between Euclidean distance and level
units.

We now remove four exceptional classes from the averaged first
moment.  On a diameter bin $t<\operatorname{diam}\omega\le2t$,
H\"older gives
\[
 \sum\bar\mu(\omega)|\omega|
 \le\left(\sum\bar\mu(\omega)\right)^{1/3}
      \left(\sum\bar\mu(\omega)|\omega|^{3/2}\right)^{2/3}.
\]
The strip tail makes this at most
$Cb(2t)^{2/3}h^{13/12}(t/h)^{5/3}e^{-ct/(3h)}$.
Summing over $t=2^js$ is negligible relative to the lower first
moment, since $s/h=(\log h)^{1/128}$.
At the remaining diameters, lengths above $\ell_+$ cost at most
\[
 \ell_+^{-1/2}Cb(s)(s/h)^{5/2}h^{7/4}
 =Cb(s)h^{13/12}(\log h)^{-59/256}.
\]
Lengths below $\ell_-$ cost at most
$Ch^{13/12}(\log h)^{-1/8}$ by the unweighted deep-path estimate.
Finally, arches of gap below $h/\log h$ cost at most
\[
 Cb(s)s^{1/3}(h/\log h)^{3/4}
 =Cb(s)h^{13/12}(\log h)^{1/384-3/4},
\]
by summing their first-moment bounds over dyadic gaps.
Each loss is $o(h^{13/12}(\log h)^{-1/12})$.

The retained first moment is therefore still at least that order.
All retained lengths are at most $\ell_+$, so their averaged mass is
at least $ch^{-1/4}(\log h)^{-7/12}$.  If bridges contribute at least
half of this, summing over $d\in[h,2h]$ already proves the theorem.
Otherwise some strip contains arches of mass
\[
 a_h\ge c h^{-1/4}(\log h)^{-C}
\]
with diameter at most $s$, lengths in $[\ell_-,\ell_+]$, and gaps
at least $h/\log h$.  Reflection and reversal let us retain one
wall and one gap direction at a fixed-factor cost.  Translate so
these arches lie above $y=0$ and have a common left endpoint $A$.

\medskip\noindent\textbf{Two exterior connectors.}
Fix an integral $R\in[R_0,2R_0]$.  Write $B$ for a retained arch's
right endpoint and $g$ for its gap.  For each retained arch we construct
bridges of height $R$ by choosing its placement between two $x$-walls
and attaching two disjoint paths below $y=0$.  There are $\asymp R$
placements.  At each placement each connector will have mass at least
$R^{-1/4}(\log h)^{-C}$, diameter $O(R)$, and length $O(R^{4/3})$.
Throughout the construction the ratio of the target scale to the
initial scale is at most a fixed power of $\log h$: indeed
$h/\log h\le g\le s$ and $R\asymp R_0$.  Thus the power loss
$(S/r)^{-C}$ in length-controlled growth becomes only logarithmic.
All constants in the geometry below are fixed before $h$ is taken
large; in particular $s/R\to0$.

Choose \(x\)-walls at separation \(R\) with the left wall at distance \(D\in[R/3,R/2]\) in \(x\)-level units from \(A\), \(O(R)\) choices with a matching lower bound. Use the \(x\)-direction with normal angle \(-\pi/6\) (versus \(\pi/2\) for \(y\)), and a divider at an \(x\)-level near the middle between the endpoints. The arch fits inside the strip far from the two walls. From \(A,B\) inside the respective sides of the divider and strictly below \(y=0\), we construct connectors to the two exterior \(x\)-walls, all centers before their target wall, of diameter \(\le C R\), length \(\le C R^{4/3}\) and mass each \(\ge R^{-1/4}(\log h)^{-C}\). On the right use the length-controlled growth directly with normals \(-y,x\), initial scale \(c_* g\) small, exact target the right wall. The divider is avoided by small initial diameter versus gap and small transverse error on continuation.

For the left connector first grow with normals \(-y,u\) where \(u\) has angle \(-5\pi/6\). Take variable exact targets at grid distances \(S\in[(2-2\delta_*)D,(2-\delta_*)D]\) along \(u\), where \(\delta_*>0\) is sufficiently small fixed, initial scale \(c_* g\). The \(u\)-ray thus makes progress to the left in \(-x\) at slope \(1/2\) and also downwards at slope \(1/2\). Choose \(c_*\) small and ray errors sufficiently small after fixing \(\delta_*\). Since \(g/R\to0\), paths stay strictly between divider and left wall, finish with remaining \(-x\) distance between \(\delta_*D/4\) and \(2\delta_*D\), well below \(y=0\) by order \(D\). There are $\asymp R$ admissible targets $S$, each giving mass at least
$R^{-1/4}(\log h)^{-C}$, so their total mass is at least
$R^{3/4}(\log h)^{-C}$. Extract the first qualifying $u$-cut in the
terminal bin, allowing the terminal wall, with the stated confinement,
length and endpoint bounds. The suffixes over all targets cost at most
$O(R^{3/4})$. Thus the distinct extracted prefixes have mass at least
$(\log h)^{-C}$. From there grow a second time with normals \(u,-x\), reaching the exact left wall (in \(-x\)), staying strictly beyond the extracted \(u\)-cut. Take initial scale a sufficiently small fixed fraction of the remaining distance in this second construction; diameters are at most a fixed constant times that distance. Hence by choosing \(\delta_*\) sufficiently small (this second-step constant is independent of it) everything fits within \(y<0\) and on the proper side of the divider. The cut is recoverable and the resulting connector mass is \(\ge R^{-1/4}(\log h)^{-C}\).

The geometry is summarized in Figure~\ref{L:fig:connectors}. The arch
stays above $y=0$, and the divider separates the two exterior paths below
that line. The recovered $u$-cut distinguishes the two left stages.
\begin{figure}[htbp]
\centering\begin{tikzpicture}[x={(1cm,0cm)},y={(0.58cm,1cm)},scale=.85,>=Stealth]
 \fill[blue!4] (-3,-2.8)--(-3,2)--(3,2)--(3,-2.8)--cycle;
 \draw[gray] (-4,0)--(4,0) node[right,black] {$y=0$};
 \draw[thick] (-3,-2.8)--(-3,2) node[above] {left $x$-wall};
 \draw[thick] (3,-2.8)--(3,2) node[above] {right $x$-wall};
 \draw[dashed,gray] (0,-2.8)--(0,1.6);
 \node[black,font=\small] at (.35,-.9) {divider};
 \draw[very thick,blue!65!black] (-.9,0) .. controls(-1,1.3) and(-.2,1.6) .. (.1,1.05)
 .. controls (.5,.65) and(.5,1.3) .. (.9,0);
 \draw[very thick,red!65!black,->] (-.9,0) ..controls(-1.1,-.4) and(-1.2,-1.7)..(-2.25,-2.1);
 \draw[very thick,red!65!black,->] (-2.25,-2.1)..controls(-2.65,-2.3) and(-2.85,-1.7)..(-3,-1.65);
 \draw[very thick,red!65!black,->] (.9,0)..controls(1.4,-.45) and(2.3,-1.8)..(3,-2.1);
 \fill(-.9,0)circle(2pt)node[above left]{$A$};
 \fill(.9,0)circle(2pt)node[above right]{$B$};
 \draw[densely dashed,gray!75] (-2.72,-1.693)--(-1.78,-2.507) node[right,black,font=\small] {$u$-cut};
 \fill(-2.25,-2.1)circle(2pt);
 \node[blue!65!black] at (-.1,1.8){retained arch};
 \node[red!65!black,align=center] at (1.2,-2.6){right connector};
 \node[red!65!black,align=right,font=\small] (leftlabel) at (-5.15,-1.6){two stages\\on the left};
 \draw[red!65!black,thin] (leftlabel.east)--(-2.55,-2.15);
\end{tikzpicture}

\caption{Schematic two-connector construction. The blue arch is retained,
while the red paths join its ends to opposite $x$-walls. The picture shows
confinement regions and the order of construction, not the microscopic
shape of a typical path. The dashed $u$-cut separates the two stages of the left connector
and makes their split recoverable.}\label{L:fig:connectors}
\end{figure}


\medskip\noindent\textbf{Recovery and the mass balance.}
Both connectors stay below the arch and on their respective sides of
the divider. Their union with the arch is therefore a strict $x$-bridge
of height $R$ and diameter at most $C_0R$, where $C_0$ is fixed.
Translate its left endpoint to the prescribed starting port. Count the input arch, the placement $D$,
and the two connector paths as labels on the resulting bridge. The
construction counts distinct connector paths by the first-cut argument;
internal growth stages are not extra labels.
If $W_R$ denotes this labelled output mass, multiplicativity at ports
and the order $R$ choices of $D$ give
\begin{equation}
 W_R\ge cR\bigl(R^{-1/4}(\log h)^{-C}\bigr)^2a_h
       \ge cR^{1/2}(\log h)^{-C}a_h.
 \label{L:eq:connector-labelled-mass}
\end{equation}

We next recover these labels from the bridge. For any strict $x$-bridge
$\omega$, let $M(\omega)$ be the number of $y$-grid lines above both
endpoints that meet $\omega$ at exactly two crossing ports. In every
constructed output the translated line $y=0$ is one of these lines:
the arch lies strictly above it and both connectors lie strictly below.
Choose this line in an output. Its two crossings determine the subarc
above the line, hence the input arch and its left and right endpoints
$A,B$. Translating $A$ back to the prescribed arch source fixes the
original placement. In particular the gap $g$ and the distance $D$
from $A$ to the left $x$-wall are determined; $R$ is already the output
height. The two remaining subarcs determine the connector paths.

The split of the left connector is also recoverable. The data $R,D,g$
specify its terminal $u$-bin and all the prefix tests used to extract
the direction-change cut. The second part stays strictly beyond that
cut in $u$, so it changes neither the earlier crossing counts nor the
prefix tests. The first qualifying cut on the completed connector is
therefore the cut selected in the construction. Here the first part
is counted as a distinct extracted prefix: its former terminal target
$S$ and the growth stages used before extraction are not retained as
labels. For the second part, its starting port and exact target wall
are now known, so its stage cuts are recovered by the same rule.
Thus each choice of the $y$-line determines at most a fixed number of
input labels, allowing for the finitely many orientation conventions.
The inverse multiplicity is at most $C M(\omega)$.

We need a bound on large values of $M$ that is uniform in the order of
the moment. Let $\mathcal G_R$ be all strict $x$-bridges from the fixed
source, of height $R$ and diameter at most $C_0R$. For
$1\le k\le c_0R$, with $c_0>0$ fixed, we claim
\begin{equation}
 \sum_{\omega\in\mathcal G_R}\mu(\omega)
       \binom{M(\omega)}{k}
 \le C^kR^2\left(C\sqrt{R/k}\right)^{k-1},
 \label{L:eq:double-cut-factorial}
\end{equation}
where $C$ is independent of $k$ and $R$.
To prove this, list the selected levels in increasing order,
$t_1<\cdots<t_k$. There are $O(R)$ choices for the far endpoint and
$O(R)$ choices for $t_1$. The two initial arms, from the two endpoints
to $t_1$, lie below that level between the $x$-walls. Each has bounded
total mass by the convex boundary-flux estimate; a further lower wall
can be placed beyond the diameter bound. Drop mutual avoidance between
these arms for the upper bound.

Between $t_i$ and $t_{i+1}$, connectivity and the two-crossing condition
give two strict $y$-bridges. If $j_i=t_{i+1}-t_i$, their total mass is
at most $B_{j_i}^2\le Cj_i^{-1/2}$. The remaining top arc has bounded
mass by the half-plane arch identity. The positive gaps satisfy
$\sum_i j_i\le C_1R$, for a fixed $C_1$ allowing for the level units.
With $a=k/R$, exponential weighting gives
\[
 \begin{aligned}
 \sum_{\substack{j_1,\ldots,j_{k-1}\ge1\\
                  \sum_i j_i\le C_1R}}
       \prod_{i=1}^{k-1}j_i^{-1/2}
 &\le e^{C_1k}
       \left(\sum_{j\ge1}j^{-1/2}e^{-kj/R}\right)^{k-1}\\
 &\le e^{C_1k}\left(C\sqrt{R/k}\right)^{k-1}.
 \end{aligned}
\]
This proves \eqref{L:eq:double-cut-factorial}, including $k=1$ with
the empty-product convention. Counting the selected levels in
increasing order is what gives $\binom Mk$.

Set $T_R=\sqrt R\log R$. Since $M(\omega)\le CR$ on
$\mathcal G_R$, for $1\le k\le T_R/2$ the preceding estimate yields
\[
 \begin{aligned}
 \sum_{\substack{\omega\in\mathcal G_R\\M(\omega)>T_R}}
       \mu(\omega)M(\omega)
 &\le \frac{CR}{\binom{\lfloor T_R\rfloor}{k}}
       \sum_{\omega\in\mathcal G_R}\mu(\omega)
       \binom{M(\omega)}{k}\\
 &\le CR^{5/2}\sqrt k
       \left(\frac{C\sqrt k}{\log R}\right)^k.
 \end{aligned}
\]
Choose $k=\lfloor c(\log R)^2\rfloor$, with $c>0$ sufficiently
small. For large $R$ this lies in both required ranges, and the last
bound is smaller than every fixed inverse power of $R$. The labelled
mass of constructed outputs with $M>T_R$ is at most a constant times
this bound, so it is negligible compared with
\eqref{L:eq:connector-labelled-mass}. On the remaining outputs the
inverse multiplicity is at most $CT_R$. Hence their unlabelled mass
is at least
\[
 \frac{cR\bigl(R^{-1/4}(\log h)^{-C}\bigr)^2a_h}
      {\sqrt R\log R}
 \ge h^{-1/4}(\log h)^{-C},
\]
after changing $C$, since $R\in[R_0,2R_0]$. This is the required
mass lower bound at each such height $R$.

The connector lengths are at most $CR^{4/3}$, uniformly over the
placements.  Since $R\le2h(\log h)^{1/64}$, their combined length is
at most $C'h^{4/3}(\log h)^{1/48}<\ell_+$ for sufficiently large $h$.
The retained arch has length between $\ell_-$ and $\ell_+$, so every
output counted above lies in $[\ell_-,2\ell_+]$.  Summing the per-height
mass over $R\in[R_0,2R_0]$ and rotating back to the fixed bridge
orientation gives
\begin{equation}
 \sum_{1\le H\le2R_0}
 \mu\{\omega\in\mathcal B_H:\ell_-\le|\omega|\le2\ell_+\}
 \ge c h^{3/4}(\log h)^{-C}.
 \label{L:eq11}
\end{equation}
The earlier bridge alternative obeys the same estimate because
$[h,2h]\subset[1,2R_0]$.  This proves the theorem.
\end{proof}

\begin{proof}[Completion of Theorem~\ref{L:thm:finite}]
Section~\ref{L:part6} proves the bridge mass estimate. The cap construction
in Section~\ref{L:part10} proves the arch lower bound with a fixed diameter
cutoff, and the gap comparison there proves the matching upper bound. Transverse confinement is the second tightness
consequence in Section~\ref{L:part10}; its first moment is
Proposition~\ref{L:prop:box-moment}. Section~\ref{L:step18} proves the
strip diameter tail. The polygon assertion is \eqref{L:eq10}, proved at
the end of Section~\ref{L:part9}.
\end{proof}
